\chapter{Система СИ, иерархия масс и Стандартная модель}
\label{ch:si-sm}

\section{Полная лестница пространственных масштабов от планковской ячейки до макромира}
\label{sec:full-ladder}

В предыдущих главах были установлены законы движения фазовой жидкости на микроуровне~$\Mlayer$,
геометрия однотактового кубооктаэдра гранецентрированной кубической решётки (\latGCKshort)
и значение постоянной тонкой структуры~$\alpha_{\mathrm{fs}}$ (\cref{ch:carrier,ch:evolution,ch:matter,ch:alpha}).
Внутреннее устройство возбуждений первых этажей --- одиночного планкона, его надпорогового ореола~$\rho_\Theta$
и трёхфокусных барионных клик~$K_3$ --- подробно исследуется в томе этажей
(\fvolrefs{ch:floors-preface,ch:floor0,ch:floor1,ch:floor2}).
Прежде чем переходить к вычислению размерных констант и спектра масс элементарных частиц,
проследим единую шкалу длин, связывающую шаг решётки $dh=\ell_P$ с атомными, молекулярными и макроскопическими масштабами.

Всякий пространственный масштаб~$\lambda$ в дискретном пространстве измеряется безразмерным числом шагов
относительно ребра однотактового кубооктаэдра $a = dh = \ell_P$ (\cref{ch:descent,def:step-tick}).
Переход от одного этажа шкалы к следующему не требует введения новых постулатов:
на каждом шаге мы рассматриваем устойчивые стационарные конфигурации того же самого автомата~$g$,
прочитанные либо на следующей координационной сфере графа решётки, либо после крупномасштабного биномиального усреднения~$\mathcal{B}$ (\cref{sec:readout,ch:macro}).

\begin{definition}[Безразмерный пространственный масштаб и комптоновский якорь]
\label{def:hop-count}
Пусть $L$ --- характерная длина физической структуры на макроуровне~$\Tlayer$, а $a = dh$ --- шаг решётки на~$\Mlayer$ (\cref{def:step-tick}); после отождествления с единицами СИ (\cref{def:si-bridge}) имеем $\ell_P := a$.
Безразмерный масштаб длины определяется отношением
\begin{equation}
  \frac{\lambda}{\ell_P}
  = \frac{L}{\ell_P},\qquad
  \ell_P := a.
\end{equation}
Инфракрасным электромагнитным якорем шкалы служит приведённая комптоновская длина волны электрона~$\lambda_C$,
отвечающая безразмерному числу шагов
\[
  N_c := \frac{\lambda_C}{\ell_P} = \frac{m_P}{\massElectron} = \frac{\caKappaGeom^{-1} m_{\arg}}{\massElectron} \approx 2{,}379\times 10^{22}.
\]
\end{definition}

Вдоль единой оси~$\lambda/\ell_P$ последовательно действуют три геометрических механизма укрупнения масштаба,
каждый из которых господствует в своей области расстояний:

\paragraph{1. Геометрический рост координационных сфер ($k=0\ldots 5$).}
В ближней зоне вокруг ядра дефекта масштаб растёт по степеням координационного числа звезды~\latGCKshort:
\[
  \lambda = \geomNeighborTwelve^{m}\,\ell_P,\qquad \geomNeighborTwelve = 12,\quad m = 0,1,2,\ldots
\]
Полоса $m=3\ldots 5$ ($\lambda/\ell_P \sim 1{,}7\times 10^3\ldots 2{,}5\times 10^5$) образует лептонное окно этажа~1 (\fvolrefs{ch:floor1}).

\paragraph{2. Восьмибитная башня электрослабого ослабления ($k=6\ldots 8$).}
Между геометрическим окном звезды и комптоновским масштабом частиц действует битовая иерархия фазового кольца
глубины $\vevHierarchySteps = \lfloor\caBitBudget\rfloor - 1 = 8$ (\cref{sec:bit-budget}):
каждый бит иерархии растягивает длину в $\alpha_{\mathrm{fs}}^{-1}\approx 137{,}036\approx \geomNeighborTwelve^2=144$ раза.
Восемь ступеней $\alpha_{\mathrm{fs}}^{-8}$ выводят на электрослабый масштаб~$v$ ($k=6$),
девятая ступень $\alpha_{\mathrm{fs}}^{-9}$ фиксирует комптоновский радиус барионного треугольника~$K_3$ протона ($N_p = m_P/m_p \approx 1{,}304\times 10^{19}$),
а десятая ступень $\alpha_{\mathrm{fs}}^{-10}$ с фазовым делителем одетого кластера $N_\varphi = 1 + \geomNeighborTwelve = 13$
приводит непосредственно к комптоновскому якорю электрона $N_c \approx 2{,}379\times 10^{22}$ ($k=8$).

\paragraph{3. Инфракрасная атомно-молекулярная лестница ($k\ge 7$).}
В окрестности комптоновского якоря~$N_c$ каждый следующий кулоновский этаж отличается множителем $\alpha_{\mathrm{fs}}^{-1}$:
\begin{equation}
\label{eq:ir-hop-ladder}
  \lambda_{k+1} = \frac{\lambda_k}{\alpha_{\mathrm{fs}}},\qquad
  \lambda_{-1} = \alpha_{\mathrm{fs}}\,\lambda_C\quad(r_e),\quad
  \lambda_0 = \lambda_C,\quad
  \lambda_1 = \frac{\lambda_C}{\alpha_{\mathrm{fs}}}\quad(a_0),\quad
  \lambda_k = \frac{\lambda_C}{\alpha_{\mathrm{fs}}^{k}}\ \ (k\ge 0).
\end{equation}
На ступени $k=9$ ($\lambda_1 = a_0$, одиннадцатая степень $\alpha_{\mathrm{fs}}^{-(\vevHierarchySteps+d)}=\alpha_{\mathrm{fs}}^{-11}$)
формируется первая боровская оболочка атома водорода с целым числом шагов $N_{a0}=\mathrm{round}(N_c/\alpha_{\mathrm{fs}})\in\mathbb{Z}$.
На ступени $k=10$ атомные ромбические додекаэдры Вигнера--Зейтца складываются во второй кристалл
с точно целым числом упаковки планковских ячеек Вороного:
\[
  N_{\mathrm{pack}}(K) = \frac{V_{\mathrm{WS}}(K\cdot dh)}{\mathrm{d}V} = \frac{\caKappaGeom\,(K\cdot dh)^3}{\caKappaGeom\,dh^3} = K^3 \in \mathbb{Z}.
\]
При этом в силу сокращения масштабного множителя~$K$ в отношении сфер $\caKappaGeom(K)=\caKappaGeom$
и в кванте действия $s_{\mathrm{floor}}(K)=(p_0/K)(K\cdot dh)=\ladderAction=\hbar/2$
все угловые каналы кубооктаэдра $(V,E,\square,\triangle)$ и параметры фазового кольца $\mathbb{Z}_{N_{\mathrm{ring}}}$
воспроизводятся на атомном и молекулярном этажах ($k=9\ldots 11$) без искажений.

\begin{proposition}[Сводная шкала длин $\lambda$ от планковской ячейки до макромира]
\label{prop:ladder-scale-steps}
Последовательность характерных масштабов при $\geomNeighborTwelve=12$ и $\alpha_{\mathrm{fs}}^{-1}\approx 137{,}036$ имеет вид:

\smallskip
\begin{center}
\footnotesize
\setlength{\tabcolsep}{4pt}
\begin{tabular}{@{}r r r l l l@{}}
\toprule
$k$ & $\lambda/\ell_P$ & $\log_{10}(\lambda/\ell_P)$ & степень $\alpha_{\mathrm{fs}}$ & физический объект & уровень \\
\midrule
0 & $1$ & $0$ & $m_P$ & планковская ячейка $\hv$ ($\lambda=\ell_P$) & $\Mlayer$ \\
1 & $12$ & $1{,}1$ & --- & каузальная звезда $\geomNeighborTwelve$; радиус одевания $R_{\mathrm{dress}}=a$ & $\Mlayer$ \\
2 & $144$ & $2{,}2$ & --- & вторая координационная сфера ($N_{\mathrm{shell}}(2)=42$) & $\Mlayer$ \\
3 & $1728$ & $3{,}2$ & --- & нижняя граница лептонного окна этажа~1 & $\Mlayer$ \\
4 & $2{,}07\times 10^4$ & $4{,}3$ & --- & средняя полоса лептонного этажа & $\Mlayer$ \\
5 & $2{,}49\times 10^5$ & $5{,}4$ & --- & верхняя граница лептонного окна ($B=0$) & $\Mlayer$ \\
6 & $4{,}95\times 10^{16}$ & $16{,}7$ & $\alpha_{\mathrm{fs}}^{-8}$ & электрослабый масштаб~$v$ и барионные клики~$K_3$ & $\Mlayer\to\Tlayer$ \\
7 & $\alpha_{\mathrm{fs}} N_c$ & $20{,}2$ & $\alpha_{\mathrm{fs}}^{-9}N_\varphi$ & классический радиус~$r_e$ и масштаб протона~$N_p$ & $\Mlayer$--$\Tlayer$ \\
8 & $N_c$ & $22{,}4$ & $\alpha_{\mathrm{fs}}^{-10}N_\varphi$ & комптоновский якорь электрона~$\lambda_C$ & $\Mlayer\to\Tlayer$ \\
9 & $N_{a0}=\mathrm{round}(N_c/\alpha_{\mathrm{fs}})$ & $24{,}5$ & $\alpha_{\mathrm{fs}}^{-11}N_\varphi$ & боровский радиус атома водорода~$a_0$ & $\Tlayer$ \\
10 & $N_c/\alpha_{\mathrm{fs}}^2$ & $26{,}6$ & $\alpha_{\mathrm{fs}}^{-12}N_\varphi$ & ковалентные связи и упаковка $N_{\mathrm{pack}}(K)=K^3\in\mathbb{Z}$ & $\Tlayer$ \\
11--13 & $10^{28}\ldots 10^{33}$ & $28{,}8\ldots 33{,}1$ & $\alpha_{\mathrm{fs}}^{-13\ldots -15}$ & биомолекулы, катушки хроматина, клетка, ткань & $\Tlayer$ \\
14 & $\sim 10^{35}$ & $35{,}3$ & $\alpha_{\mathrm{fs}}^{-16}$ & макроскопический метр, гидродинамика тела & $\Tlayer$ \\
15 & $\lambda\gg\sigma_R a$ & --- & --- & непрерывный предел огрубления (теор.~\ref{thm:t-cr}) & $\Tlayer$ \\
\bottomrule
\end{tabular}
\end{center}
\end{proposition}

\begin{remark}
Подчеркнём различие между внутренним радиусом одевания ядра электрона $R_{\mathrm{dress}}=a$ на первой сфере кубооктаэдра ($k=1$)
и его инфракрасной комптоновской длиной волны $\lambda_C = N_c \ell_P$ ($k=8$):
первый характеризует размер неразложимой топологической сердцевины на~$\Mlayer$,
тогда как вторая определяет масштаб стоячей фазовой огибающей на~$\Tlayer$, формирующейся после десяти ступеней ослабления~$\alpha_{\mathrm{fs}}$.
\end{remark}

\section{Переход от естественных единиц ячейки к системе СИ}

На микроуровне~$\Mlayer$ все механические величины выражаются через шаг решётки $a = dh$, такт времени $\htick = dt$
и первичный квант фазового действия $\ladderAction = \planckHbar\,\Delta\varphi_{\min} = \planckHbar/2$ (\cref{ch:descent,ch:matter}).
Связь этих внутренних единиц с лабораторной системой СИ полностью определяется геометрическим коэффициентом однотактового кубооктаэдра $\caKappaGeom = R_{\mathrm{in}}/R_{\mathrm{out}} = 1/\sqrt{2}$ (\cref{ch:carrier}).

\begin{definition}[Планковская сшивка с системой СИ]
\label{def:si-bridge}
Положим шаг решётки равным планковской длине $\ell_P := a$. Тогда с учётом коэффициента вписанной сферы кубооктаэдра~$\caKappaGeom$ основные масштабы микроуровня~$\Mlayer$ и макроуровня~$\Tlayer$ связаны соотношениями:
\begin{align}
  \htick &= \caKappaGeom\, t_P, \\
  \ladderEnergy &= \caKappaGeom\, E_P, \\
  c &= \caKappaGeom\, \caSignalSpeed.
\end{align}
\end{definition}

В классической литературе планковское время вводят как $t_P = \ell_P / c$, неявно отождествляя скорость переноса по ребру решётки~$\caSignalSpeed$ с макроскопической скоростью света~$c$ (что отвечало бы $\caKappaGeom=1$).
На гранецентрированной кубической решётке изотропный сферический волновой фронт вписан в однотактовый кубооктаэдр с отношением радиусов $\caKappaGeom = 1/\sqrt{2}$,
поэтому истинный такт автомата короче $t_P$ в $\caKappaGeom^{-1}=\sqrt{2}$ раз ($\htick = \caKappaGeom t_P$),
а скорость звена равна $\caSignalSpeed = \ell_P/\htick = c/\caKappaGeom = \sqrt{2}\,c$ (\cref{sec:two-c}).

В базовых размерностях длины~$\mathrm{L}$, массы~$\mathrm{M}$ и времени~$\mathrm{T}$ планковские величины имеют вид:
\begin{align}
  [t_P] &= \mathrm{T}, &
  [\ell_P] &= [c\,t_P] = \mathrm{L}, \\
  [m_P] &= \left[\frac{\planckHbar}{c^{2}t_P}\right] = \mathrm{M}.
  \label{eq:si-dim-planck-base}
\end{align}
Безразмерный геометрический множитель~$\caKappaGeom$ сохраняет размерности времени и скорости:
\begin{align}
  [\htick] &= \mathrm{T}, &
  [\caSignalSpeed] &= \frac{\ell_P}{\htick} = \mathrm{L\,T^{-1}}, \\
  [c] &= [\caKappaGeom\,\caSignalSpeed] = \mathrm{L\,T^{-1}}.
  \label{eq:si-dim-speed}
\end{align}
Отсюда вся механическая лестница первичных квантов ячейки на паре $(\ell_P,\htick)$ однозначно воспроизводит размерности величин в системе СИ:
\begin{align}
  [\ladderAction] &= [\planckHbar] = \mathrm{M\,L^{2}T^{-1}}, \\
  [\ladderEnergy] &= \left[\frac{\ladderAction}{\htick}\right] = \mathrm{M\,L^{2}T^{-2}}, \\
  [\ladderMomentum] &= \left[\frac{\ladderAction}{\ell_P}\right] = \mathrm{M\,L\,T^{-1}}, \\
  [\ladderForce] &= \left[\frac{\ladderMomentum}{\htick}\right] = \mathrm{M\,L\,T^{-2}}, \\
  [u_P] &= \left[\frac{m_P c^{2}}{\ell_P^{3}}\right] = \mathrm{M\,L^{-1}T^{-2}}
  \quad\text{($\NotationUnitJoulePerCubicMetre$, планковская плотность энергии).}
  \label{eq:si-dim-ladder}
\end{align}
При этом постоянная тонкой структуры~$\alpha_{\mathrm{fs}}$ с самого начала является чистым безразмерным числом решётки ($[\alpha_{\mathrm{fs}}]=1$, \cref{ch:alpha}), а электрический заряд~$e_0$ служит лишь переводом целого топологического инварианта вихря $n\in\mathbb{Z}$ в кулоны.

\section{Постоянная тонкой структуры}

Полный вывод постоянной тонкой структуры~$\alpha_{\mathrm{fs}}$ из геометрии однотактового кубооктаэдра ($\caKappaGeom = R_{\mathrm{in}}/R_{\mathrm{out}}$, мягкая грань $(1-\caKappaGeom^6)/\caKappaGeom^6 = 7$) и кулоновского знаменателя $\hat M = 1 + \geomNeighborTwelve(\lfloor\caBitBudget\rfloor - 1) = 97$ изложен в главе~\ref{ch:alpha} (\cref{thm:alpha,def:alpha-meaning}).
В настоящей главе это безразмерное число используется как уже установленная константа носителя, управляющая электрослабым, барионным и лептонным масштабами масс.

\section{Электрослабый масштаб и вакуумное среднее}

Из информационного бюджета Бекенштейна одной планковской ячейки $\caBitBudget = 2\pi/\ln 2 \approx 9{,}065$ (\cref{sec:bit-budget})
следует разрядность фазового регистра $\mathrm{mod\_bits} = \lfloor\caBitBudget\rfloor = 9$ ($N_{\mathrm{ring}} = 2^9 = 512$).
За вычетом одного старшего бита занятости ядра планкона ($b\in\{0,1\}$) остаётся ровно
\[
  \vevHierarchySteps = \lfloor\caBitBudget\rfloor - 1 = 8
\]
иерархических битовых ступеней ослабления амплитуды между первичным квантом ячейки~$\ladderEnergy$ и электрослабой пеной макроуровня~$\Tlayer$.

\begin{theorem}[Электрослабый масштаб вакуумного среднего]
\label{thm:vev}
В естественных единицах ячейки $(dh,dt)$, где $\ladderEnergy = \ladderAction/dt = \caKappaGeom E_P$, а фазовый вычет одного такта равен $r = \Delta\varphi_{\min}/(2\pi) = 1/(4\pi)$, электрослабый масштаб вакуумного среднего равен
\[
  v = \frac{1}{\sqrt{r}}\,\alpha_{\mathrm{fs}}^{\vevHierarchySteps}\,\ladderEnergy
  = \alpha_{\mathrm{fs}}^{8} E_P \sqrt{2\pi}
  \approx 246{,}09\ \mathrm{ГэВ}.
\]
\end{theorem}

\begin{proof}
\pfromref{информационного бюджета ячейки}{sec:bit-budget}
каждая из $\vevHierarchySteps = \lfloor\caBitBudget\rfloor - 1 = 8$ битовых ступеней иерархии вносит кулоновский фактор ослабления~$\alpha_{\mathrm{fs}}$ (\cref{thm:alpha}), что даёт пиковую шкалу энергии
\[
  v_{\mathrm{peak}} = \alpha_{\mathrm{fs}}^{\vevHierarchySteps} E_P = \frac{\alpha_{\mathrm{fs}}^{\vevHierarchySteps}}{\caKappaGeom}\,\ladderEnergy.
\]
При переходе от узловой амплитуды на~$\Mlayer$ к макроскопической гауссовой огибающей на~$\Tlayer$ (\cref{thm:t-cr,def:soft-readout})
нормировочный интеграл по фазовому циклу с вычетом $r = \Delta\varphi_{\min}/(2\pi)$ и геометрическим фактором кубооктаэдра $\caKappaGeom^2 = 1/2$ даёт множитель $(\caKappaGeom\sqrt{r})^{-1} = \sqrt{2\pi}$.
\psubst{$v_{\mathrm{peak}}$}
получаем требуемое выражение $v = r^{-1/2}\alpha_{\mathrm{fs}}^{\vevHierarchySteps}\ladderEnergy \approx 246{,}09\ \mathrm{ГэВ}$.
\end{proof}

\section{Спектр масс частиц и магнитные моменты нуклонов}

Начиная от первичного кванта инерции ячейки $m_{\arg} = \ladderEnergy/c^2 = 2\ladderAction dt/dh^2 = \caKappaGeom m_P$ (\cref{sec:true-mass-carrier}),
массы бозона Хиггса, нуклонов, кварковых ступеней и трёх поколений заряженных лептонов однозначно определяются геометрией кубооктаэдра~\latGCKshort\ без подгоночных констант Юкавы (\fvolrefs{ch:floor2}):

\begin{proposition}[Иерархия масс и магнитных моментов в единицах $m_{\arg}$]
\begin{enumerate}[leftmargin=*]
\item \textbf{Бозон Хиггса и топ-кварк на полке~$v$ ($\alpha_{\mathrm{fs}}^{\vevHierarchySteps}$):}
  бозон Хиггса отвечает квадрату отношения сфер кубооктаэдра~$\caKappaGeom^2$, а верхний потолок одиночного фермионного фокуса (топ-кварк) --- первой степени~$\caKappaGeom$:
  \begin{align*}
    m_H^{(\mathrm{bare})} &= \caKappaGeom^2 v = \frac{\caKappaGeom^2}{\sqrt{r}}\,\alpha_{\mathrm{fs}}^{\vevHierarchySteps}\,m_{\arg} \approx 123{,}04\ \mathrm{ГэВ},\qquad
    m_H = m_H^{(\mathrm{bare})}\sqrt{1 + \vevHierarchySteps\frac{\alpha_{\mathrm{fs}}}{\pi\caKappaGeom^2}} \approx 125{,}31\ \mathrm{ГэВ},\\
    m_t^{(\mathrm{bare})} &= \caKappaGeom\,v = \frac{m_H^{(\mathrm{bare})}}{\caKappaGeom} \approx 174{,}01\ \mathrm{ГэВ}.
  \end{align*}
\item \textbf{Барионный треугольник~$K_3$: протон, нейтрон и их магнитные моменты ($\alpha_{\mathrm{fs}}^{\vevHierarchySteps+1}$):}
  с учётом кванта упаковки кубооктаэдра $\delta_{\mathrm{pack}} = \caKappaGeom^2/\geomNeighborTwelve = 1/24$:
  \begin{align*}
    m_p &= \alpha_{\mathrm{fs}}\,m_H^{(\mathrm{bare})}\bigl(1 + \delta_{\mathrm{pack}}\bigr) \approx 935{,}30\ \mathrm{МэВ},\qquad
    m_n = m_p + 2\massElectron \approx 936{,}33\ \mathrm{МэВ},\\
    \mu_p &= \bigl(d - 5\delta_{\mathrm{pack}}\bigr)\mu_N = \frac{67}{24}\,\mu_N \approx +2{,}7917\,\mu_N,\qquad
    \mu_n = \bigl(-2 + 2\delta_{\mathrm{pack}}\bigr)\mu_N = -\frac{23}{12}\,\mu_N \approx -1{,}9167\,\mu_N,\\
    \mu_p + \mu_n &= \bigl(1 - d\,\delta_{\mathrm{pack}}\bigr)\mu_N = \frac{7}{8}\,\mu_N = 0{,}8750\,\mu_N.
  \end{align*}
\item \textbf{Легчайший заряженный дефект (электрон, $\alpha_{\mathrm{fs}}^{\vevHierarchySteps+2}/N_\varphi$) и три поколения лептонов ($e,\mu,\tau$):}
  одетый кластер из $N_\varphi = 1 + \geomNeighborTwelve = 13$ ячеек определяет массу электрона
  \[
    \massElectron = \frac{\alpha_{\mathrm{fs}}^2\,m_H}{N_\varphi} \approx 0{,}5133\ \mathrm{МэВ},
  \]
  а три пространственных стека поколений ($g\in\{0,1,2\}$, $N_{\mathrm{gen}}=d=3$) с амплитудой кубооктаэдра~$\caKappaGeom^{-1}$
  и начальным фазовым сдвигом двухкомпонентного спинора на битовом регистре ячейки $\delta_0 = 2/\lfloor\caBitBudget\rfloor = 2/9$
  дают формулу масс поколений:
  \[
    \sqrt{m_g} \propto 1 + \frac{1}{\caKappaGeom}\cos\!\left(\frac{2\pi g}{d} + \frac{2}{\lfloor\caBitBudget\rfloor}\right),\qquad
    \frac{m_\mu}{\massElectron} \approx 206{,}7703\ \ (\text{PDG: }206{,}7683),\qquad
    \frac{m_\tau}{\massElectron} \approx 3477{,}47\ \ (\text{PDG: }3477{,}23\pm 0{,}23).
  \]
\item \textbf{Кварковые ступени ($u,d,s,c,b$):}
  зарядовые шаги переключения вершин треугольника~$K_3$ равны $m_u = 2^2\massElectron = 4\massElectron$, $m_d = d^2\massElectron = 9\massElectron$ ($m_u/m_d = 4/9$),
  а старшие кварковые пороги равны $m_s = m_p/(\geomNeighborTwelve-2)$, $m_c = \frac{d+1}{d}m_p$ и $m_b = \frac{1-\caKappaGeom^6}{d\,\caKappaGeom^6}\alpha_{\mathrm{fs}}v$.
\end{enumerate}
\end{proposition}

\section{Сводка численных сверок с данными СИ и PDG}
\label{sec:si-pdg-experiments-journal}

Все приведённые выше соотношения являются прямыми алгебраическими следствиями аксиом носителя и не содержат подгоночных параметров.
Численное сопоставление теоретических значений с эталонами CODATA и таблицами Particle Data Group (PDG)
оформлено в серии критических опытов \textbf{CE-M} (\cref{sec:lab-series-ce-m,ch:critical-construction,ch:computational-laboratory})
по метрологическому протоколу главы~\ref{ch:measurement-processing}.

% Auto-generated by scripts/gen_lab_records_tex.py — do not edit.
% Generated: 2026-11-16 10:34 UTC
% Journal prose: book/sources/lab-journal/ (hand-authored).
\begin{table}[htbp]
\centering
\caption{Серия CE-M: модель (upstream §8.2) vs эталон PDG/CODATA. Не натурный эксперимент --- воспроизводимый прогон реестра.}
\label{tab:si-pdg-ce-m}
\small
\begin{tabular}{@{}lrrrl@{}}
\toprule
Наблюдаемая & Модель & Эталон & $|\delta|$ & Verify \\
\midrule
$v$ [GeV] & 246{,}09 & 246{,}22 & $5{,}39\times 10^{-4}$ & \texttt{Higgs\_mass} \\
$m_H$ [GeV] & 125{,}31 & 125{,}25 & $4{,}73\times 10^{-4}$ & \texttt{Higgs\_mass} \\
$m_p$ [GeV] & 0{,}9353 & 0{,}9383 & $3{,}16\times 10^{-3}$ & \texttt{Proton\_mass} \\
$m_e$ [GeV] & 0{,}000513 & 0{,}000511 & $4{,}50\times 10^{-3}$ & \texttt{Electron\_mass} \\
$m_n$ [GeV] & 0{,}9363 & 0{,}9396 & $3{,}44\times 10^{-3}$ & \texttt{Neutron\_mass} \\
$\alpha^{-1}$ & 137{,}035999 & 137{,}035999 & $7{,}19\times 10^{-11}$ & \texttt{Alpha\_si\_bridge} \\
$\sin^2\theta_W$ (bare) & 0{,}23077 & 0{,}23122 & $1{,}95\times 10^{-3}$ & \texttt{Alpha\_bridges} \\
$\alpha_s(M_Z)$ & 0{,}11788 & 0{,}11790 & $1{,}36\times 10^{-4}$ & \texttt{---} \\
\bottomrule
\end{tabular}
\end{table}


\noindent
Подробные постановки критических опытов CE-M-00 ($\caSignalSpeed$ и~$c$), CE-M-01 (макроскопический световой конус), CE-M-02/03 (константа~$\alpha_{\mathrm{fs}}$ и спектр масс) и CE-M-04 (полное квантование возбуждений на~$\Mlayer$) приведены в главе~\ref{ch:critical-construction} и приложении~\ref{app:si-pdg-experiments}.

\section{Единая природа четырёх фундаментальных взаимодействий}
\label{sec:four-forces-map}

На микроуровне~$\Mlayer$ действует единственный локальный закон эволюции~$g$ (\cref{ch:evolution}) на решётке~\latGCKshort.
Четыре фундаментальных взаимодействия классической и квантовой физики не являются четырьмя независимыми полями, наложенными поверх пространства:
на макроуровне~$\Tlayer$ они возникают как четыре различных канала проявления одной и той же фазовой упругости кипящего океана в зависимости от топологии дефекта и масштаба расстояний.

\begin{proposition}[Электромагнитный канал]
\label{prop:channel-em}
На макроуровне~$\Tlayer$ градиенты общей фазы~$U(1)$ на рёбрах кубооктаэдра и циркуляция потока~$\Phi_\square$ по граням ячейки после биномиального усреднения~$\mathcal{B}$ порождают уравнения Максвелла (\cref{ch:macro,def:soft-readout}).
Скорость распространения поперечных волн фазового давления равна $c = \sqrt{K_P/\mu_P} = \caKappaGeom\caSignalSpeed$,
а безразмерная константа связи кулоновского взаимодействия двух топологических вихрей $Q = n e_0$ на больших расстояниях равна~$\alpha_{\mathrm{fs}}$ (\cref{ch:alpha,thm:alpha}).
\end{proposition}

\begin{proposition}[Слабый канал и угол Вайнберга из геометрии одетого кластера]
\label{prop:weinberg-bare}
Слабое взаимодействие отвечает переключению относительной фазы двух компонент спинора~$\mathbb{C}^2$ ($SU(2)$-переворот на полуоборот кольца $N_{\mathrm{ring}}/2 = 256$ тактов, \cref{ch:evolution}).
Отношение числа пространственных измерений $d=3$ к полному числу ячеек одетого однотактового кластера $N_\varphi = 1 + \geomNeighborTwelve = 13$ (центральное ядро и двенадцать соседей кубооктаэдра) задаёт затравочное значение квадрата синуса угла смешивания Вайнберга:
\begin{equation}
\label{eq:sin2-bare}
  \sin^2\theta_W^{(\mathrm{bare})}
  = \frac{d}{\geomNeighborTwelve + 1}
  = \frac{d}{N_\varphi}
  = \frac{3}{13} \approx 0{,}23077.
\end{equation}
\end{proposition}

\begin{proof}
Из $N_\varphi = 13$ независимых фазовых степеней свободы одетого кластера ровно $d=3$ пространственных направления участвуют в векторном переносе импульса по осям решётки, тогда как остальные $N_\varphi - d = 10$ принадлежат чисто внутренним циркуляциям оболочки.
Отношение пространственной доли к полному фазовому бюджету кластера даёт~\eqref{eq:sin2-bare}, а учёт одного электромагнитного шага по $2\vevHierarchySteps = 16$ спин-иерархическим состояниям башни определяет значение на полюсе $Z$-бозона:
\[
  \sin^2\theta_W(M_Z) = \frac{d}{N_\varphi} + \frac{\caKappaGeom^2\alpha_{\mathrm{fs}}}{\vevHierarchySteps} = \frac{3}{13} + \frac{\alpha_{\mathrm{fs}}}{16} \approx 0{,}231225
  \qquad\bigl(\text{PDG: }0{,}23122\pm 0{,}00004,\ 0{,}13\sigma\bigr).
\]
\end{proof}

\begin{proposition}[Сильный канал и константа связи $\alpha_s(v)$]
\label{prop:alpha-s-seed}
На малых расстояниях внутри барионного треугольника~$K_3$ три взаимно ортогональные координатные плоскости $\{YZ,XZ,XY\}$ его рёбер образуют трёхцветный замок удержания ($N_c = d = 3$, \fvolrefs{ch:floor2}).
На электрослабой полке~$v$ (глубина иерархии $\vevHierarchySteps = 8$) константа сильного взаимодействия определяется отношением числа цветовых плоскостей $d=3$ к полному фазовому объёму восьмибитной башни $\vevHierarchySteps\pi$:
\begin{equation}
\label{eq:alpha-s-seed}
  \alpha_s(v) = \frac{d}{\vevHierarchySteps\,\pi} = \frac{3}{8\pi} \approx 0{,}11937,
\end{equation}
что при стандартном однопетлевом сдвиге от $v\approx 246{,}1\ \mathrm{ГэВ}$ к полюсу $M_Z\approx 91{,}19\ \mathrm{ГэВ}$ согласуется с измерением PDG ($\alpha_s(M_Z) = 0{,}1179\pm 0{,}0009$).
\end{proposition}

\begin{proof}
В то время как электрослабое взаимодействие проходит все $\vevHierarchySteps = 8$ битовых ступеней ослабления по степеням~$\alpha_{\mathrm{fs}}$, короткодействующий барионный канал замыкает три ортогональные плоскости клики~$K_3$ непосредственно внутри одной иерархической ячейки башни объёма~$\vevHierarchySteps\pi$, откуда следует~\eqref{eq:alpha-s-seed}.
\end{proof}

\begin{theorem}[Три сектора взаимодействий как сечения кубооктаэдра: высота~$\vevHierarchySteps$, ширина~$N_\varphi$ и площадь~$\hat M$]
\label{thm:three-interaction-sectors}
Три калибровочных взаимодействия Стандартной модели не требуют слияния в единую численную константу~$\alpha_{\mathrm{GUT}}$ на сверхвысокой энергии, поскольку они представляют собой три взаимно ортогональных геометрических сечения одного и того же шага~$g$ двухкомпонентного спинора $z\in\mathbb{Z}_{N_{\mathrm{ring}}}[i]^2$ на однотактовом кластере кубооктаэдра:
\begin{enumerate}[noitemsep]
\item \textbf{Сильный сектор (высота иерархической башни $\vevHierarchySteps = 8$):}
  короткодействующее замыкание барионного треугольника~$K_3$ по $d=3$ ортогональным плоскостям внутри одной ячейки башни без разведения по $12$ лучам звезды даёт $\alpha_s(v) = d/(\vevHierarchySteps\pi) = 3/(8\pi)$.
\item \textbf{Слабый сектор (ширина однотактового кластера $N_\varphi = 1 + \geomNeighborTwelve = 13$):}
  переключение компонент спинора $z_1\leftrightarrow z_2$ по $d=3$ пространственным осям внутри $13$-узлового кластера даёт $\sin^2\theta_W^{(\mathrm{bare})} = d/N_\varphi = 3/13$.
\item \textbf{Электромагнитный сектор (полная площадь «звезда $\times$ башня» $\hat M = 1 + (N_\varphi - 1)\vevHierarchySteps = 97$):}
  дальнодействующий перенос общей фазы~$U(1)$ по всем $\geomNeighborTwelve = N_\varphi - 1 = 12$ лучам звезды на всех $\vevHierarchySteps = 8$ этажах башни плюс центральное ядро даёт $\alpha_0 = \caKappaGeom/\hat M \to \alpha_{\mathrm{fs}}$.
\end{enumerate}
\end{theorem}

\begin{remark}[Гравитационный канал как макроскопическая упругость океана]
\label{rem:channel-gr-open-in-book}
Гравитация не является отдельным квантовым полем на~$\Mlayer$: присутствие массивных вихревых узлов (запирающих фазовое натяжение~$\mathrm{Arg}$ в ячейках с жёсткостью $K_P \approx 4{,}63\times 10^{113}\ \mathrm{Па}$) создаёт дальнодействующую деформацию плотности фазовых шагов окружающего океана,
которая после усреднения~$\mathcal{B}$ описывается эффективной метрикой ОТО с гравитационной постоянной $G = \ell_P c^2 / m_P = \caKappaGeom\,dh\,c^2 / m_{\arg}$.
\end{remark}

\begin{proposition}[Полные матрицы смешивания кварков (CKM) и нейтрино (PMNS) из $N_\varphi=13$ секторов и гауссовой факторизации $\hat M=97$]
\label{prop:ckm-pmns-sectors}
\begin{enumerate}[leftmargin=*]
\item \textbf{Гауссова факторизация $\hat M = 97$ в $\mathbb{Z}[i]$ и матрица CKM ($\lambda,A,\bar\rho,\bar\eta,\gamma$):}
  в гауссовом кольце спинора~$\mathbb{Z}[i]$ кулоновский знаменатель раскладывается в сумму квадратов внутреннего ветвления квадратной грани $d_{\mathrm{in}}(\square)=4$ и битового бюджета ячейки $\lfloor\caBitBudget\rfloor=9$:
  \[
    \hat M = d_{\mathrm{in}}(\square)^2 + \lfloor\caBitBudget\rfloor^2 = 4^2 + 9^2 = 97.
  \]
  На удвоенном по спину кластере $2N_\varphi = 26$ комплексная координата вершины треугольника унитарности кварков равна
  \[
    \bar\rho + i\,\bar\eta = \frac{d_{\mathrm{in}}(\square) + i\,\lfloor\caBitBudget\rfloor}{2\,N_\varphi} = \frac{4 + 9i}{26}
    \quad\Longrightarrow\quad
    \bar\rho = \frac{2}{13}\approx 0{,}1538\ (\text{PDG: }0{,}159\pm 0{,}010),\quad
    \bar\eta = \frac{9}{26}\approx 0{,}3462\ (\text{PDG: }0{,}348\pm 0{,}010),
  \]
  откуда длина стороны треугольника равна $R_b = |\bar\rho+i\bar\eta| = \sqrt{\hat M}/(2N_\varphi) = \sqrt{97}/26 \approx 0{,}3788$ (PDG: $0{,}3826\pm 0{,}011$),
  а угол CP-нарушения связан с лептонной фазой $\delta_0 = 2/\lfloor\caBitBudget\rfloor = 2/9$ тождеством
  \[
    \tan\gamma = \frac{\bar\eta}{\bar\rho} = \frac{\lfloor\caBitBudget\rfloor}{d_{\mathrm{in}}(\square)} = \frac{1}{2\delta_0} = \frac{9}{4}
    \quad\Longrightarrow\quad
    \gamma = \arctan\!\left(\frac{9}{4}\right) \approx 66{,}04^\circ
    \qquad\bigl(\text{PDG: }65{,}4^\circ\pm 2{,}5^\circ,\ 0{,}25\sigma\bigr).
  \]
  Параметры Вольфенштейна $\lambda$ и~$A$ равны
  \[
    \lambda = \frac{d}{N_\varphi}\left(1 - \frac{d-1}{d}\,\delta_{\mathrm{pack}}\right) = \frac{35}{156} \approx 0{,}22436\ (\text{PDG: }0{,}22500\pm 0{,}00067),\qquad
    A = \frac{\geomNeighborTwelve - 2}{\geomNeighborTwelve} = \frac{5}{6} \approx 0{,}8333\ (\text{PDG: }0{,}826\pm 0{,}015).
  \]
\item \textbf{Разбиение $N_\varphi = 13$ секторов вершинным каналом $V=(1,4,7)$ и углы смешивания нейтрино PMNS:}
  разложение $N_\varphi = 1_{\mathrm{core}} + d_{\mathrm{in}}(V) + d_{\mathrm{same}}(V) + d_{\mathrm{out}}(V) = 1 + 1 + 4 + 7 = 13$ однозначно фиксирует все три угла смешивания нейтральных лептонов:
  \[
    \sin^2\theta_{12}^{\mathrm{PMNS}} = \frac{d_{\mathrm{same}}(V)}{N_\varphi} = \frac{4}{13} \approx 0{,}3077\ (\text{PDG: }0{,}307\pm 0{,}013,\ 0{,}05\sigma),
  \]
  \[
    \sin^2\theta_{23}^{\mathrm{PMNS}} = \frac{d_{\mathrm{out}}(V)}{N_\varphi} = \frac{7}{13} \approx 0{,}5385\ (\text{PDG: }0{,}546\pm 0{,}021,\ 0{,}36\sigma),
  \]
  \[
    \sin^2\theta_{13}^{\mathrm{PMNS}} = \frac{1 + d_{\mathrm{in}}(V)}{N_{\mathrm{shell}}(3)} = \frac{2}{92} = \frac{1}{46} \approx 0{,}02174\ (\text{PDG: }0{,}0220\pm 0{,}0007,\ 0{,}37\sigma).
  \]
\end{enumerate}
\end{proposition}
